\section{Error analysis, robustness, uncertainty, and limitations}
\label{sec:errors}

\subsection{Q2e: Validation error analysis}
The finalised Q2 candidate (best CV Macro-F1, see \texttt{results/models/q2\_best.json}) was analysed on
validation only. All misclassified validation rows are in \texttt{results/tables/q2\_errors\_all.csv}.
\inputtable{q2_errors}
\inputtable{q2_error_groups}
\inputtable{q2_error_numeric_profile}

% TODO(student): interpret
\studenttodo{Q2e: categorise at least ten errors into patterns (e.g.\ confident FPs with
\texttt{poutcome}=success, FNs among first-time contacts with no history, seasonal months, \texttt{contact}
missing); connect them to input characteristics; propose concrete improvements (features, representation,
threshold, model class) without using test labels.}

\subsection{Q5a: Cross-validation variability}
\inputtable{q5_repeated_cv}
\fig[0.85]{q5_repeated_cv_box}{Fold-wise Macro-F1 of frozen candidates.}

\subsection{Q5b: Robustness}
\inputtable{q5_robustness}
\inputtable{q5_robustness_retention}
\fig{q5_robustness}{Validation Macro-F1 vs.\ corruption severity (mean $\pm$ SD over 5 repetitions).}

% TODO(student): interpret
\studenttodo{Q5a/b: which models are stable across folds; which degrade fastest under noise/missingness and
why (distance-based kNN vs.\ trees' thresholds vs.\ linear models; how imputation interacts with missingness).
Is the scenario realistic for this domain?}

\subsection{Q5c: Paired bootstrap}
\inputtable{q5_bootstrap}
\fig[0.7]{q5_bootstrap}{Bootstrap distribution of the Macro-F1 difference.}

% TODO(student): interpret
\studenttodo{Q5c: what the interval permits (uncertainty of the difference on this validation distribution;
whether 0 is excluded) and does not permit (no causal/general claim, does not cover training randomness or
other datasets, not a test-set statement, multiple comparisons). State sampling unit and $B$.}

\subsection{Limitations and threats to validity}
% TODO(student): interpret
\studenttodo{Threats to validity of the full study: single dataset per task, campaign data from 2008--2010
(temporal shift), no temporal split, subjective interpretability/implementation scores and weights, limited
grids, validation reused for several analyses, Macro-F1 at a fixed 0.5 threshold, etc.}
